\section{Additional benchmark results}
\label{app:grid}

\subsection{The full parameter grid}

Table~\ref{tab:grid} gives the best-over-implementations sparse latency divided
by the best-over-implementations dense latency, across head width and across
batch, at each agent count and at the most favourable sparsity available. Head
width is swept to $512$ and batch from $32$ to $1024$. The largest batch we
attempted, $4096$ at $512$ agents, could not be measured: the score matrix alone
runs to several gigabytes and the device is shared, so that configuration is
absent.

\begin{table}[ht]
\centering\small
\caption{Best sparse divided by best dense, over head width (upper block) and
batch (lower block). Values below one would indicate sparsity winning.}
\label{tab:grid}
\begin{tabular}{rrrr}
\toprule
head width $d$ & $N{=}96$ & $N{=}192$ & $N{=}512$ \\
\midrule
16 & 0.66 & 0.71 & 0.74 \\
32 & 0.51 & 0.69 & 0.62 \\
64 & 0.87 & 1.32 & 2.15 \\
128 & 1.38 & 2.17 & 1.88 \\
256 & 2.10 & 2.07 & 2.11 \\
512 & 1.66 & 1.72 & 2.18 \\
\bottomrule
\end{tabular}

\vspace{1em}

\begin{tabular}{rrrr}
\toprule
batch $B$ & $N{=}96$ & $N{=}192$ & $N{=}512$ \\
\midrule
32 & 0.55 & 0.65 & 0.57 \\
1024 & 0.55 & 0.75 & 0.72 \\
\bottomrule
\end{tabular}%

\end{table}

\subsection{Dispersion behind the headline numbers}
\label{app:dispersion}

The primary statistic is the minimum over interleaved samples, on the modelling
assumption that contention can only add time. That statistic is trustworthy only
if the underlying sample was not wildly dispersed, so the harness records the
median, the interquartile range and the spread of per-run medians for every
configuration, and the released artifacts carry all of them. Two properties of
this host limit what any of it can mean. Both GPUs carried between 30 and
80~GiB
of other tenants' allocations and near-continuous utilisation throughout the
campaign, and clocks could not be locked. We therefore report the spread of
per-run medians beside every configuration rather than a confidence interval: a
ratio of minima has no sampling distribution.

\subsection{Drift control}
\label{app:drift}

The headline sweep was repeated in full on the same device, hours later and
under different load from other tenants. Table~\ref{tab:drift} gives how far
each agent count moved between the two runs, in absolute latency and in the
ratio against the best dense implementation, which is the quantity every claim
is stated in.

\begin{table}[ht]
\centering\small
\caption{Movement between two runs of the identical sweep. From $192$ agents
upward both the latencies and the ratios reproduce to about one per cent. At
$96$ they do not, by a wide margin, which is why no claim in this paper is read
off an agent count below $192$.}
\label{tab:drift}
\begin{tabular}{rrrrrr}
\toprule
& & \multicolumn{2}{c}{\textbf{latency} $|\Delta|$ (\%)}
& \multicolumn{2}{c}{\textbf{ratio} $|\Delta|$ (\%)} \\
\cmidrule(lr){3-4}\cmidrule(lr){5-6}
$N$ & configurations & median & max & median & max \\
\midrule
96 & 44 & 0.1 & 4.2 & 0.3 & 2.5 \\
192 & 44 & 0.2 & 6.1 & 0.2 & 6.0 \\
384 & 44 & 0.7 & 11.3 & 0.5 & 11.6 \\
512 & 44 & 6.1 & 12.0 & 7.2 & 12.6 \\
\bottomrule
\end{tabular}%

\end{table}

From $192$ agents upward the measurement is stable: median movement is under
one per cent in latency and under half a per cent in the ratio, with the worst
single configuration under eleven per cent. At $96$ agents it is not stable at
all --- every implementation ran between $14\%$ and $394\%$ slower in the second
run, and the fastest configurations in the first run had recorded inter-run
spreads above $85\%$ at the time. That is the signature of a launch-bound
regime on a contended device rather than of anything about attention, and it is
why the crossover of Section~\ref{sec:fused} is reported only from $192$ agents
upward even though the sweep measures smaller ones.

The repeat is on the same physical card, so it establishes stability across
runs and hours, not independence from one device, and we were not able to
obtain a second card free of other tenants. It covers the sweep of
Section~\ref{sec:composed} and Section~\ref{sec:fused}'s figure; the tiling
search of Table~\ref{tab:autotune} was run once, so the crossover it reports
carries the reproducibility of a single session rather than of a repeat.
Measurements were taken on both cards of this host --- the core sweeps, the
search and the stage accounting on one, the head-width sweep and the
pre-score measurements on the other --- and no cross-device calibration is
claimed, which is a further reason no figure from one is divided by a figure
from the other.

\section{Implementation and measurement validation}
\label{app:verification}

\subsection{Functional equivalence}

Every implementation is checked against the reference at every configuration
whose latency any released artifact reports: $158$ combinations of agent count,
head width and retained-peer count, $136$ of them with $k < N-1$ so that
selection is actually exercised, over $1.25$ million query rows in total. The
check matters more than it looks. Exact top-$k$ is discontinuous in the scores,
so two implementations that accumulate the score product in different orders can
rank a near-tied pair differently and aggregate over different peers, producing
a large output difference from a rounding-level score difference. An earlier
version of our per-query fused kernel computed the score row as a reduction
instead of through the matrix-multiply units and disagreed with the reference at
a handful of configurations for exactly that reason; the kernels reported here
use the same product as the reference.

With that fixed, every implementation agrees with the reference to
\texttt{fp32} rounding --- maximum deviation $3.2\times10^{-6}$, and zero for
the two implementations that share the reference's operator sequence exactly ---
with one attributable exception. The fused kernel that selects by threshold
retains every peer at or above the $k$-th largest score, so an exact tie at that
boundary retains $k{+}1$ peers where a rank-based selection retains $k$. Four
query rows in $1.25$ million have such a tie, at three of the $158$
configurations, and those four rows are the only place any implementation
differs from the reference by more than rounding. The check is run by
\texttt{scripts/verify\_all\_configurations.py}, which classifies a tie
divergence separately from a failure so that the distinction is in the artifact
and not only in this paragraph.

\subsection{What each implementation executes}

The masked and gathered implementations compute the same function while
executing different work, which is the distinction the paper turns on, so
regression tests pin it. Maximum absolute deviation
between the two paths is below $10^{-10}$ in \texttt{float64}. A FLOP counter
attached to the running program reports identical totals for dense and masked
--- the counter scores matrix products, and masked issues the same two as dense
--- against a smaller total for gathered, the difference being exactly the
aggregation term. Tracing every matrix product confirms that the gathered path
issues no dense aggregation product and does issue the compacted one, and that
every top-$k$ variant still issues the full score product, which is assumption
A2 made mechanical.

\subsection{Which backend serves the fused dense baseline}

A probe forces each scaled-dot-product backend in turn, at both dtypes and at
three agent counts, under no mask, an additive float mask and a Boolean mask,
and the profiler records which kernel the default dispatcher launches. The
FlashAttention backend is rejected under any non-null mask in either form, and
in \texttt{fp32} under any input, so the library dense baseline runs the
memory-efficient CUTLASS kernel. Forcing FlashAttention on the same unmasked
\texttt{bf16} shape is faster than the kernel we do get, so a baseline able to
reach it would widen rather than narrow the gap we report --- which is part of
why the dense baseline in this revision is the better of the library path and a
kernel we wrote.

\subsection{The host allocator control}
\label{app:cpu}

Dense, masked and gathered each allocate and free several multi-MiB score
matrices per forward. The default allocator returns those pages to the kernel
and they are re-faulted on the next call, which at $N{=}192$ adds roughly
$37{,}000$ minor page faults per forward and inflates latency by up to
$8.8\times$. The tax is proportional to the bytes of large temporaries each
implementation churns, so it differs per implementation and does not cancel in a
ratio; in an un-pinned run it inverted the ordering at one configuration. We
raise the allocator's retention threshold so that freed temporaries are kept,
and count minor faults per forward from the process's own resource usage across
the timed region. The harness records the allocator environment, the thread
counts, the affinity mask and the fault totals in each run's metadata, and the
superseded un-pinned run is retained beside the pinned one.

\subsection{Machine balance}
\label{app:roofline}

Peak arithmetic rate is measured with a large matrix multiply and peak bandwidth
with a large copy, in the same process as the measurements they normalise. That
matters here more than it usually would: on this shared host the peak we measure
varies between $53$ and $77$~TFLOP/s across sessions, depending on what else is
running. Every fraction-of-peak figure in the paper is against the peak measured
in the same process as the number it normalises, and is therefore an upper bound
on the true efficiency rather than an estimate of it. Comparisons between two
implementations measured in one process are unaffected; comparisons of a
fraction across sections are not, and we make none. The
copy source is written before it is read: an untouched allocation on this
operating system is backed by the shared zero page, so a host-side copy from it
re-reads one cache line and reports a bandwidth several times the machine's
real figure, and an earlier version of this probe did exactly that: every CPU
ridge point derived from it was too large by roughly a factor of six. Device
work is also timed with device events rather than host wall-clock, since host
timing around a sub-millisecond copy measures launch and synchronisation. The
probe reports two ridge points per device, one with a cache-resident buffer and
one with a buffer larger than the last-level cache, because the two differ by a
factor of two on the host and a claim that holds for one need not hold for the
other.

\section{Additional training analysis}
\label{app:training}

\subsection{Pre-registration chronology}
\label{app:prereg}

The analysis plan was committed when eight of the forty runs in the
pre-registered family had produced a final evaluation and before any of their
returns had been inspected. Commit history establishes the ordering: the
positive-control artifacts at 12:18, the plan at 12:56, and the remaining
artifacts at 19:11 the same day. The plan has two commits, the second appending
to its deviations log; the endpoint, comparisons, test, effect sizes, interval
and correction are byte-identical to the pre-unblinding version. Timestamps
establish ordering and nothing more. Data generation had demonstrably already
begun.

One correction is logged against the released plan. It asserts that at most one
comparison in the family can reach corrected significance. That is false: Holm
is a step-down procedure, so once a first rejection is obtained its later
thresholds are looser and further rejections remain possible. The manuscript
states the corrected version; the deviation is logged against the released plan.

\subsection{The remaining comparisons}

Five further comparisons were registered and all are Holm-adjusted nulls. Two
deserve a note. The comparison at twelve agents has the one sign opposite to
its registered direction and carries the family's largest point estimate, and
we set it aside on the ground that neither variant learns the task there. That
ground was not registered: it is a post-hoc exclusion of the largest observed
effect, we have logged it as a deviation against the released plan, and a
reader who rejects it should read the comparison as an under-powered null with
a large point estimate, which is what the registered analysis returns either
way. The evidence for the exclusion is this. A uniform-random policy scores
$-758.9$ under a forty-episode evaluation,
against $-864.9$ for the dense variant, which is worse, and $-717.9$ for the
sparse variant, which is better but barely learning; the comparison is between a
diverged run and a struggling one. Separately, for two comparisons the bootstrap
interval excludes zero while the permutation test does not. The registration
named the permutation test, so it governs; the disagreement is unsurprising at
five seeds, where the interval is coarse, discrete and of uncertain coverage. We
report both so that the estimated effect and the strength of evidence for it are
both visible.

Descriptively, and supporting no claim: the budgeted variant meets its budget
but not by the intended route, its gate collapsing to an expected support near
one within the first tenth of training while the dual multiplier decays to zero;
the entmax variant settles at a support of $1.64$ of five peers; and attention
entropy falls from $0.35$ nats over five peers to $0.12$ over two. The three sit
on different supports and are reported only as an instrumentation check.

\subsection{The partially observable environment}
\label{app:env}

The variant modifies the standard cooperative task in one respect: a peer's
relative-position entry is zeroed when the peer is farther than a fixed radius.
Reward, action space, landmark observations, episode length and observation
dimension are unchanged. Two properties limit what it can support. The mask is a
zeroing rather than an explicit visibility flag, so a peer at relative position
exactly zero is not distinguishable from an absent one; measured under a uniform
random policy, this affects at most $0.15\%$ of visible peers. And the visible
fraction falls with the number of agents while the absolute number visible
rises, because the world size is fixed, which is a difficulty confound that
rules out across-agent-count return comparisons.

The communicating variants carry between $7.2\%$ and $19.0\%$ more trainable
parameters than the communication-free baseline, since the attention module adds
parameters the baseline has no analogue for. The positive controls therefore
establish that the communicating architecture outperforms the baseline, not that
the channel alone accounts for the difference.

\section{Literature audit}
\label{app:audit}

For each audited paper we read the full text, including appendices and the
official supplement where one exists --- through the ar5iv rendering of the
arXiv version where the paper is on arXiv, and the venue's own copy
otherwise --- searched it for seven temporal keywords,
and recorded whether it makes a computational-efficiency or complexity claim
about its mechanism, whether it reports any measured temporal quantity
attributable to that mechanism, and the figure, table or sentence that settled
the decision. A statement such as ``training took a few days'' does not qualify,
since it is not attributable to the mechanism. Scope matters for at least one
entry, whose affirmative verdict rests on a wall-clock figure in an appendix and
would be negative on the main body alone. This is an audit of reporting practice
in a small, deliberately selected set --- the papers we ourselves relied on ---
and not a sample of either literature; a negative verdict means the argument
rests on operation counts or asymptotics, not that it is wrong.

\begin{table}[ht]
\centering\small
\caption{The twelve audited papers, their efficiency claim, the kind of evidence
it rests on, and whether any measured temporal quantity attributable to the
mechanism appears anywhere in the paper or its supplement. A ``no'' means the
argument rests on operation counts or asymptotics, not that it is wrong.}
\label{tab:audit}
\begin{tabular}{p{0.29\linewidth}p{0.36\linewidth}p{0.14\linewidth}l}
\toprule
\textbf{Paper} & \textbf{Efficiency claim} & \textbf{Type} &
\begin{tabular}[b]{@{}l@{}}\textbf{Measured}\\\textbf{time?}\end{tabular} \\
\midrule
\multicolumn{4}{l}{\emph{Multi-agent communication and attention}} \\
\citet{das2019tarmac} & Qualitative bandwidth argument only & qualitative & no \\
\citet{jiang2018atoc} & Communication is 'restricted by bandwidth\,\ldots & agent/message counts & no \\
\citet{jiang2018dgn} & receiving a large amount of information\,\ldots & asymptotic / qualitative & no \\
\citet{singh2018ic3net} & continuous communication allows efficient\,\ldots & message counts / convergence epochs & no \\
\citet{sukhbaatar2016commnet} & none explicit & --- & n/a \\
\citet{wen2022mat} & Sequence modelling 'reduces the complexity\,\ldots & asymptotic & \textbf{yes} \\
\multicolumn{4}{l}{\emph{Efficient attention and sparse mixture-of-experts}} \\
\citet{beltagy2020longformer} & Sliding-window + global attention scales\,\ldots & asymptotic & \textbf{yes} \\
\citet{choromanski2021performers} & FAVOR+ gives linear rather than quadratic\,\ldots & asymptotic & \textbf{yes} \\
\citet{fedus2022switch} & Sparse routing increases parameters at\,\ldots & FLOPs-matched & \textbf{yes} \\
\citet{kitaev2020reformer} & LSH attention reduces O(L$2$) to O(L log L) & asymptotic & \textbf{yes} \\
\citet{shazeer2017moe} & 1000x capacity increase at minor\,\ldots & FLOPs + measured & \textbf{yes} \\
\citet{zaheer2020bigbird} & reduces this quadratic dependency to linear';\,\ldots & asymptotic + memory & no \\
\bottomrule
\end{tabular}%

\end{table}

\section{Reproduction and provenance}
\label{app:repro}

The measurement campaign is training-free and completes in under an hour on one
device. \texttt{scripts/run\_revision\_campaign.sh} regenerates every artifact
in the order the manuscript reports them;
\texttt{scripts/make\_revision\_tables.py} and
\texttt{scripts/make\_revision\_figures.py} regenerate every table and figure
from those artifacts, so no tabular number in the manuscript is transcribed by
hand; and \texttt{scripts/check\_revision\_numbers.py} re-derives the numbers
that appear in prose and fails on any disagreement.

Three provenance caveats. The hardware description and the aggregate provenance
record were written before this project's environment was upgraded and describe
a superseded software stack and a host that is no longer available; both have
been corrected, and the run-level provenance for the training variants lives in
the per-run configuration files rather than in the aggregate record. The GPUs
used for every measurement in this paper are shared with other tenants, so
absolute latencies are not reproducible to better than the drift reported in
Appendix~\ref{app:drift} and only ratios between implementations timed in the
same interleaved sweep should be compared. And the activation-memory measurement
is reported from a run with the card otherwise free, since a shared-host repeat
exhausts memory one agent count earlier.
